% !TEX root = ../general/main.tex
\section{Introduction}

Precision medicine uses heterogeneous treatment effects (HTE) to guide clinical decisions.
Researchers estimate the HTE from either a randomised controlled trial (RCT) or observational real-world data (RWD).
Randomised trials are the gold standard for causal effect estimation, but they are underpowered for subgroup effects and often have limited follow-up.
Real-world data from registries and cohorts are larger and often follow patients for longer.
Regulators increasingly accept them as evidence.
Unmeasured confounding can bias treatment effect estimates from observational data.
We combine the two data sources. We identify the effect from the trial and borrow power and longer follow-up from the real-world data.
Our motivating example is HIV antiretroviral therapy.
The treatment effect varies with baseline CD4 count, viral load, and age.
We fuse the ACTG~175 trial \citep{hammer1996} with the Multicenter AIDS Cohort Study \citep{kaslow1987}.
The cohort followed participants over more than 25 years with biannual visits.

The fusion of a trial with real-world data raises three challenges for survival outcomes.
The first is unmeasured confounding in the real-world data.
Treatment is not randomised, and assignment depends on factors the measured covariates do not fully capture.
The no-unmeasured-confounders assumption is untestable and often violated in practice \citep{vanderweele2017evalue}.
The second is heterogeneity between the two data sources.
Even when the treatment effect transports across sources, the baseline prognosis often does not.
The trial and observational populations typically differ in eligibility, follow-up, and case mix.
A naive analysis confounds this baseline difference with the treatment effect.
This heterogeneity is not confined to the mean: differences in measurement protocol and follow-up quality also reshape the error distribution across sources.
The third is the survival outcome.
Right censoring arises when the event has not occurred by the end of follow-up, so the survival time is known only to exceed the censoring time.
The real-world data may also be interval-censored, with events recorded only at periodic visits.

Unmeasured confounding can be addressed through a confounding function.
The confounding function measures the confounding bias in the real-world data.
It was introduced as a sensitivity-analysis device within a single study \citep{robins2000sensitivity}.
The confounding function can be estimated by combining a randomised and an observational study.
\citet{Kallus2018} learn the confounding function from the RCT and debias the observational treatment effect.
\citet{YangLiuZengWang2025} develop a semiparametric-efficient estimator that jointly estimates the treatment effect and the confounding function, and derive a test for unmeasured confounding.
The elastic integrative estimator of \citet{yang2023elastic} uses that test to decide whether to borrow from the real-world data, and discards it when confounding is detected.
Two frequentist methods extend the approach to survival outcomes.
\citet{Ye2025} work under an accelerated failure time model with right censoring and high-dimensional covariates, and treat the confounding function as a sparse component selected jointly with the treatment effect.
\citet{mao2025statistical} target the conditional restricted mean survival time under right censoring.


Bayesian methods combine data sources by borrowing information through the prior.
The power prior of \citet{ibrahim2000power} raises the external likelihood to a fractional power that controls the borrowing.
The meta-analytic-predictive prior of \citet{neuenschwander2010} makes the borrowing hierarchical and adaptive to between-source heterogeneity.
Later variants tie the borrowing to the between-source discrepancy through commensurate priors \citep{hobbs2012commensurate}.
These methods all target a marginal or control-arm effect.
More recent literature targets heterogeneous effects directly.
\citet{zhou2021incorporating} pool trial and external outcome surfaces with Bayesian additive regression trees, but do not allow for unmeasured confounding in the external source.
\citet{dimitriou2025causal} develop a multi-task Gaussian process with a data-adaptive borrowing parameter.
None of these Bayesian methods targets heterogeneous treatment effects on survival outcomes under unmeasured confounding.

Bayesian additive regression trees \citep{chipman2010bart} are a popular method for causal machine learning \citep{hill2011bayesian}.
They are supported by theoretical guarantees \citep{rockova2020posterior} and have strong empirical performance in both estimation accuracy and uncertainty quantification \citep{dorie2019, thal2023, kabata2026}.
The uncertainty quantification is intrinsic to the model and follows directly from the posterior distribution.
Black-box predictors such as (deep) neural networks are flexible alternatives, but do not provide intrinsic uncertainty quantification and do not separate the causal effect from confounding bias.
We will show that a structured approach tailored to the problem at hand and its statistical constraints outperforms a generic black-box competitor in both accuracy and calibration.

We develop a Bayesian nonparametric framework to estimate heterogeneous treatment effects on survival outcomes.
We target the acceleration factor, a causally interpretable estimand on the survival-time scale.
We fuse a randomised trial with real-world data without assuming the observational source is unconfounded.
A confounding function absorbs the bias from unmeasured confounding in the real-world data.
We accommodate right- and interval-censored data, and exploit the longer follow-up of the real-world source for robustness.
A hierarchical Dirichlet process mixture models the error distribution nonparametrically and lets it vary across sources.